\#{ resumeow — "sleek": refined two-column sidebar layout. Same data + same render.py as default.tex; only this file differs. The main column is a flowfram text flow, so competences/experiences of any length paginate by themselves: page 1 sits beside the full sidebar, later pages beside a thin accent strip. }
% !TeX program = pdflatex
\documentclass[11pt, a4paper]{article}

\usepackage[T1]{fontenc}
\usepackage[utf8]{inputenc}
\usepackage[british]{babel}
\usepackage[left = 0mm, right = 0mm, top = 0mm, bottom = 0mm]{geometry}
\usepackage[stretch = 25, shrink = 25, tracking = true, letterspace = 40]{microtype}
\usepackage{graphicx}
\usepackage{xcolor}
\usepackage{marvosym}
\usepackage{enumitem}
\setlist{parsep = 0pt, topsep = 0pt, partopsep = 1pt, itemsep = 1pt, leftmargin = 5mm}

\usepackage{FiraSans}
\renewcommand{\familydefault}{\sfdefault}

\definecolor{accent}{HTML}{\VAR{ color }}
\definecolor{sidebar}{HTML}{\VAR{ color }}
\definecolor{toolbg}{HTML}{EEF0F4}

% Inline "tool chip" -- the LaTeX equivalent of a markdown `backtick` span.
% Used directly in the YAML data, e.g. \tool{GitLab CI}.
\newcommand{\tool}[1]{{\setlength{\fboxsep}{2pt}\colorbox{toolbg}{\footnotesize\ttfamily\textcolor{accent}{#1}}}}

\newcommand{\dates}[1]{\hfill\mbox{\small\itshape #1}}
\newcommand{\smaller}[1]{{\footnotesize #1}}
\newcommand{\evensmaller}[1]{{\scriptsize #1}}
% Sidebar header: white small caps + faint rule.
\newcommand{\headleft}[1]{\vspace*{2.2ex}{\footnotesize\bfseries\scshape\textls[50]{#1}}\par%
    \vspace*{-1.2ex}{\color{white!55!sidebar}\hrulefill}\par\vspace*{0.8ex}}
% Main header: accent small caps + short accent underline.
\newcommand{\headright}[1]{\vspace*{2.6ex}{\Large\bfseries\scshape\color{accent}\textls[50]{#1}}\par%
    \vspace*{0.1ex}{\color{accent}\rule{3.2em}{1.2pt}}\par\vspace*{1.0ex}}
% A key/transversal competence: bold accent title, body underneath. No box
% around it, so a long list flows across frames; the \nopagebreak pair only
% forbids a break *between* the title and its body.
\newcommand{\competence}[2]{{\bfseries\color{accent}#1}\par\nopagebreak\vspace{0.3em}\nopagebreak#2\par\vspace{0.7em}}

%%%%%%%%%%%%%% PAGE FRAMES %%%%%%%%%%%%%%%%%%%%%%%%%%%%%%%%%%%%%%%%%%%%%%%%%%
% Every frame below is derived from these lengths: tune them, not the frame
% definitions.
\newlength{\cvband}\setlength{\cvband}{0.33\paperwidth}     % page-1 sidebar
\newlength{\cvstrip}\setlength{\cvstrip}{0.1\paperwidth}    % page 2+ accent strip
\newlength{\cvpad}\setlength{\cvpad}{6.5mm}                 % sidebar inner padding
\newlength{\cvgutter}\setlength{\cvgutter}{8mm}             % band -> text gap
\newlength{\cvrmargin}\setlength{\cvrmargin}{14mm}
\newlength{\cvtmargin}\setlength{\cvtmargin}{10mm}
\newlength{\cvbmargin}\setlength{\cvbmargin}{12mm}

\usepackage{flowfram}

% Coloured bands: full-bleed, no content of their own.
\newstaticframe[1]{\cvband}{\paperheight}{0pt}{0pt}[band]
\newstaticframe[>1]{\cvstrip}{\paperheight}{0pt}{0pt}[strip]
\setstaticframe*{band,strip}{backcolor=sidebar}

% The sidebar text sits in its own frame, inset inside the band (page 1 only).
\newstaticframe[1]{\dimexpr\cvband-2\cvpad\relax}%
  {\dimexpr\paperheight-\cvtmargin-\cvbmargin\relax}{\cvpad}{\cvbmargin}[sbtext]
\setstaticframe*{sbtext}{textcolor=white}

% The main column: one continuous flow. Narrow beside the sidebar on page 1,
% wide beside the strip afterwards -- LaTeX moves the text between them.
\newflowframe[1]{\dimexpr\paperwidth-\cvband-\cvgutter-\cvrmargin\relax}%
  {\dimexpr\paperheight-\cvtmargin-\cvbmargin\relax}%
  {\dimexpr\cvband+\cvgutter\relax}{\cvbmargin}[maincol]
\newflowframe[>1]{\dimexpr\paperwidth-\cvstrip-\cvgutter-\cvrmargin\relax}%
  {\dimexpr\paperheight-\cvtmargin-\cvbmargin\relax}%
  {\dimexpr\cvstrip+\cvgutter\relax}{\cvbmargin}[maincolwide]

\usepackage[colorlinks = true, urlcolor = white, linkcolor = accent]{hyperref}

\begin{document}
\setlength{\parindent}{0pt}
\setlength{\parskip}{0pt}
\setlength{\fboxsep}{0pt}
\pagestyle{empty}
\raggedbottom

%%%%%%%%%%%%%% SIDEBAR CONTENT (static, page 1) %%%%%%%%%%%%%%%%%%%%%%%%%%%%%
\setstaticcontents*{sbtext}{%
\raggedright
\null\hfill\includegraphics[width=0.75\linewidth]{\VAR{ img }}\hfill\null
\vspace*{1.5ex}
\begin{center}{\Large\bfseries\textls[30]{\VAR{ name }}}\end{center}

\headleft{\VAR{ i18n.profile }}
\BLOCK{ if profile is string }
\smaller{\VAR{ profile }}
\BLOCK{ else }
\begin{itemize}
\BLOCK{ for item in profile }
\item \smaller{\VAR{ item }}
\BLOCK{ endfor }
\end{itemize}
\BLOCK{ endif }

\headleft{\VAR{ i18n.desiderata }}
\BLOCK{ if desiderata is string }
\smaller{\VAR{ desiderata }}
\BLOCK{ else }
\begin{itemize}
\BLOCK{ for item in desiderata }
\item \smaller{\VAR{ item }}
\BLOCK{ endfor }
\end{itemize}
\BLOCK{ endif }

\headleft{\VAR{ i18n.languages }}
\begin{itemize}
\BLOCK{ for lang in languages }
\item \smaller{\VAR{ lang }}
\BLOCK{ endfor }
\end{itemize}

\headleft{\VAR{ i18n.contacts }}
\smaller{%
\MVAt\ \VAR{ contacts.mail } \\[0.5ex]
\Mobilefone\ \VAR{ contacts.phone }
\BLOCK{ for link in contacts.links or [] }
\\[0.5ex]\Mundus\ \href{\VAR{ link.link }}{\VAR{ link.display }}
\BLOCK{ endfor }
}%
}

\headright{\VAR{ i18n.key_competences }}
\BLOCK{ for c in key_competences }
\competence{\VAR{ c.name }}{\smaller{\VAR{ c.text }}}
\BLOCK{ endfor }

\BLOCK{ if transversal_competences }
\headright{\VAR{ i18n.transversal_competences }}
\BLOCK{ for c in transversal_competences }
\competence{\small \VAR{ c.name }}{\evensmaller{\VAR{ c.text }}}
\BLOCK{ endfor }
\BLOCK{ endif }

\headright{\VAR{ i18n.work_experiences }}
\BLOCK{ for e in work_experiences }
{\bfseries \VAR{ e.title }}\ \textit{— \VAR{ e.company }}\dates{\VAR{ e.dates }}\par\nopagebreak\vspace{0.3em}\nopagebreak
\evensmaller{\VAR{ e.description }}
\BLOCK{ if 'items' in e }
\begin{itemize}
\BLOCK{ for it in e['items'] }
\item \evensmaller{\textbf{\VAR{ it.project }:} \VAR{ it.text }}
\BLOCK{ endfor }
\end{itemize}
\BLOCK{ endif }
\BLOCK{ if not loop.last }
\par\vspace{0.5em}{\color{accent!25}\hrulefill}\par\vspace{0.5em}
\BLOCK{ else }
\par\vspace{0.5em}
\BLOCK{ endif }
\BLOCK{ endfor }

\headright{\VAR{ i18n.education }}
\begin{itemize}
\BLOCK{ for ed in education }
\item {\bfseries \VAR{ ed.title }}\ \textit{— \VAR{ ed.institution }}\dates{\VAR{ i18n.grade_label } \VAR{ ed.grade }}
\BLOCK{ endfor }
\end{itemize}

\end{document}
